\section{Pair mass below a fixed scale}\label{sec:pairs}

We next consider a finite probability law without any separation assumption.
Its exact entropy may greatly exceed its entropy at a prescribed scale.
The following lemma detects that excess through the mass that can be
placed in fair pairs at smaller scales.  An upper bound for a sumset
will make the lemma useful for the block laws in the next section.
Varj\'u's scale-local decomposition into fair two-point submeasures
\cite[Proposition~21]{VarjuAbsolute2019} is a predecessor of this approach.
Here the sumset term controls the total capacity over all finer bands;
we prove the needed estimate by randomized nested partitions.

\Needspace{5\baselineskip}
\begin{definition}[Pair capacity]\label{def:pair-capacity}
Let $\nu$ be a probability law with finite support $F\subset\R$.
For $\ell\in\mathbb Z$, let $\mathcal E_\ell(F)$ be the set of
unordered pairs $\{x,y\}\subset F$ satisfying
$2^{-\ell}\le |x-y|<2^{1-\ell}$.
Define $c_\ell(\nu)$ as the maximum of
\[
 \sum_{e\in\mathcal E_\ell(F)}w_e
 \quad\text{subject to}\quad
 w_e\ge0,\qquad
 \frac12\sum_{e\ni x}w_e\le\nu(\{x\})\quad(x\in F).
\]
When $\mathcal E_\ell(F)$ is empty, this maximum is zero.
\end{definition}

Thus $c_\ell(\nu)$ is the largest total weight of fair two-point laws
in the indicated distance band whose weighted sum is a submeasure of
$\nu$.  Summing the constraints gives $\sum_e w_e\le1$; the feasible
set is compact, so the maximum exists and $0\le c_\ell(\nu)\le1$.
For example, if $\nu=p\delta_x+(1-p)\delta_y$ with $x\ne y$, then
the capacity in the band containing $|x-y|$ is $2\min\{p,1-p\}$,
and all its other capacities vanish.

\begin{lemma}[Finite-law pair bound]\label{lem:pair-bound}
For every integer $k\ge1$ there is a finite constant $C_k$ such that
every finitely supported probability law $\nu$ on $\R$ and every
$\rho>0$ satisfy
\begin{equation}\label{eq:pair-bound}
 k\bigl(H(\nu)-H^0_\rho(\nu)\bigr)
 \le \log|kF|
   +C_k\sum_{\ell:\,2^{-\ell}\le2\rho}c_\ell(\nu),
 \qquad F=\supp\nu,
\end{equation}
where $kF=\{x_1+\cdots+x_k:x_i\in F\}$.
The constant is independent of the smallest positive distance in $F$.
\end{lemma}

\begin{proof}
Let $X_1,\ldots,X_k$ be independent with law $\nu$, write
$\mathbf X=(X_1,\ldots,X_k)$, and put $Y=X_1+\cdots+X_k$.
Knowing the initial length-$\rho$ cells leaves
$k(H(\nu)-H^0_\rho(\nu))$ bits of uncertainty about $\mathbf X$.
Revealing $Y$ costs at most $\log|kF|$ bits.  We bound the remaining
uncertainty by successively refining those cells.  The refinement cuts
are random, independent of the samples, to prevent atoms close to a
fixed boundary from being charged too often.

\paragraph{The partitions.}
Set $\rho_h=\rho4^{-h}$ for integers $h\ge0$, and start with the
half-open grid intervals of length $\rho$.  Suppose a parent at level
$h$ has length $L\in[\rho_h/2,2\rho_h]$.  Let $N$ be a nearest
integer to $L/\rho_{h+1}$, with ties resolved by a fixed rule.
Then $2\le N\le8$.  Divide the parent into $N$ equal pieces, and
perturb its internal cuts independently and uniformly by amounts in
$[-\rho_{h+1}/10,\rho_{h+1}/10]$.
The unperturbed child lengths belong to
$[0.75\rho_{h+1},1.25\rho_{h+1}]$; after perturbation they belong
to $[0.55\rho_{h+1},1.45\rho_{h+1}]$.  Thus the cuts remain ordered
and every child has length in $[\rho_{h+1}/2,2\rho_{h+1}]$.
Conditional on the previous partitions, every new cut has density
at most
\begin{equation}\label{eq:pair-cut-density}
 \frac{5}{\rho_{h+1}}=\frac{20}{\rho_h}.
\end{equation}
All cuts are sampled independently of $\mathbf X$; a parent that
contains several coordinates uses the same cuts for each coordinate.

If $F$ is a singleton, the lemma is immediate.  Otherwise choose
an integer $J$ such that
\[
 2\rho_J<\min\{|x-y|:x,y\in F,\ x\ne y\}.
\]
Every level-$J$ cell then contains at most one point of $F$, for
every realization of the partitions.  No bound on $J$ will be needed.

For a fixed partition scheme, denote by $\mathbf C_h$ the tuple of
level-$h$ cells containing the samples.  Nesting and separation at
level $J$ give
\begin{align}
 k\bigl(H(\nu)-H^0_\rho(\nu)\bigr)
 &=H(\mathbf X\mid\mathbf C_0)\notag\\
 &\le\log|kF|+H(\mathbf X\mid\mathbf C_0,Y)\notag\\
 &=\log|kF|+
   \sum_{h=0}^{J-1}H(\mathbf C_{h+1}\mid\mathbf C_h,Y).
 \label{eq:pair-chain}
\end{align}
The whole partition scheme is sampled independently of $\mathbf X$.
We now average this finite chain over the schemes; equivalently, the
conditional entropies on its right also condition on that scheme.
For the summand at level $h$, the cell labels involved use only cuts
through level $h+1$.
Conditional on those cuts, all later cuts are independent of
$\mathbf X$ and hence of $Y$ and these labels, so they may be
discarded from the conditioning.

\paragraph{One refinement after observing the sum.}
Fix the partitions through level $h$, and condition first on a tuple
of current cells $A_1,\ldots,A_k$ of positive probability.
The coordinates remain independent, with laws $\nu(\cdot\mid A_i)$.
Write $m_i=\E X_i$ in this conditional law.  Each mean belongs to
its half-open parent, since it is an average of finitely many points
of that parent.  Fix the new cuts and let
\[
 a=\min_{\substack{1\le i\le k\\
              c\text{ a new internal cut of }A_i}}|m_i-c|.
\]
This minimum is positive almost surely in the new cuts.
Call coordinate $i$ exceptional when $|X_i-m_i|\ge a/(4k)$,
let $p_i$ be its probability, and put $t=\sum_i p_i$.
Let $E$ be the event that at least two coordinates are exceptional,
and write $q=\P(E)$.  Independence, before conditioning on $Y$, gives
\begin{equation}\label{eq:pair-two-exceptions}
 q\le\sum_{i<j}p_ip_j\le t^2,
 \qquad q\le t.
\end{equation}
Also put
\[
 B=\left\{\left|Y-\sum_i m_i\right|\ge a/2\right\}.
\]
This second event is known once $Y$ is known, and $\P(B)\le t$:
if every coordinate is nonexceptional, the sum deviation is less
than $a/4$.

On $E^c\cap B^c$, every child label equals the child label of its
mean.  Indeed, a changed label forces $|X_i-m_i|\ge a$ for some
$i$, including a change at a half-open endpoint.  There is then at
most one exceptional coordinate, so
\[
 \left|\sum_j(X_j-m_j)\right|
 \ge a-\frac{k-1}{4k}a>a/2,
\]
a contradiction.  Each coordinate has at most eight child labels.
In the next display, conditioning on the fixed partitions and new cuts
is implicit.  By revealing the flag $1_E$ in addition to $Y$, we obtain
\begin{align}
 H(\mathbf C_{h+1}\mid Y,\mathbf C_h=(A_i))
 &\le h_2(q)+3k\bigl(q+\P(B)\bigr)\notag\\
 &\le (3+6k)t\le9k\sum_i p_i.
 \label{eq:pair-flag}
\end{align}
Here $h_2(q)\le3\sqrt q$ and
\eqref{eq:pair-two-exceptions} were used.  The elementary binary
entropy bound follows from
$-(1-q)\ln(1-q)\le q$ and $-\sqrt q\ln q\le2/e$.
There is no binary-entropy cost for $B$, since it is determined by
the conditioned observation $Y$.  This distinction is what makes
\eqref{eq:pair-flag} linear in the exceptional probabilities.

\paragraph{Averaging the cuts and converting to pairs.}
For a fixed value of $X_i$, set $u=|X_i-m_i|$.
The event $a\le4ku$ requires one of at most $7k$ cuts to be within
$4ku$ of its corresponding mean.  The density bound
\eqref{eq:pair-cut-density} and a union bound give
\[
 \P_{\rm cuts}(a\le4ku)
 \le7k\frac{20}{\rho_h}(8ku)
 =1120k^2\frac{u}{\rho_h}.
\]
Consequently
\[
 \E_{\rm cuts}p_i
 \le\frac{1120k^2}{\rho_h}\E|X_i-m_i|
 \le\frac{1120k^2}{\rho_h}\E|X_i-X_i'|,
\]
where $X_i'$ is an independent copy in the same current cell.
The last inequality is Jensen's inequality.  These expectations
still use the current-cell conditional laws.

For a fixed level-$h$ partition, integrating the cells out gives
the symmetric coupling
\[
 \pi_h=\sum_{A:\,\nu(A)>0}\nu(A)
 \bigl(\nu(\cdot\mid A)\otimes\nu(\cdot\mid A)\bigr).
\]
Both its marginals equal $\nu$, and its pairs have distance at most
$2\rho_h$.  Each of the $k$ coordinate contributions becomes the
same integral against $\pi_h$.  Thus \eqref{eq:pair-flag}, averaged
over the new cuts and current cells, is at most
\begin{equation}\label{eq:pair-distance-cost}
 \frac{10080k^4}{\rho_h}\int|x-y|\,d\pi_h(x,y).
\end{equation}

For each band $2^{-\ell}\le|x-y|<2^{1-\ell}$, the corresponding
portion of $\pi_h$ supplies a feasible mixture in
Definition~\ref{def:pair-capacity}.  Explicitly, the weight assigned
to the unordered pair $\{x,y\}$ is
$\pi_h(x,y)+\pi_h(y,x)=2\pi_h(x,y)$.
The resulting submeasure has mass
$\sum_{y:\,2^{-\ell}\le|x-y|<2^{1-\ell}}\pi_h(x,y)$ at $x$,
which is at most $\nu(\{x\})$.  Its total weight equals the mass of
that band under $\pi_h$.  Therefore \eqref{eq:pair-distance-cost}
is at most
\begin{equation}\label{eq:pair-one-level}
 20160k^4
 \sum_{\ell:\,2^{-\ell}\le2\rho_h}
 \frac{2^{-\ell}}{\rho_h}c_\ell(\nu).
\end{equation}
This estimate also holds after averaging the previous partitions;
its right side refers only to the original law $\nu$.

\paragraph{Summing all finer scales.}
For a fixed $\ell$, the factors $2^{-\ell}/\rho_h$ grow by four
at each step and are summed only while they are at most two.
Hence
\[
 \sum_{h\ge0:\,2^{-\ell}\le2\rho_h}
       \frac{2^{-\ell}}{\rho_h}\le\frac83.
\]
Summing \eqref{eq:pair-one-level} in \eqref{eq:pair-chain} proves
\eqref{eq:pair-bound}; one may take $C_k=60000k^4$.
The separating depth $J$ disappears from the bound.  In particular,
the estimate includes separations smaller than any prescribed
exponential scale.
\end{proof}

For a fixed finite law only finitely many pair capacities are
nonzero.  Later we will apply Lemma~\ref{lem:pair-bound} to a finite
family of finite laws for each fixed block count.  Their nonzero
capacities therefore lie in a common finite range, although the
upper endpoint need not have any effective bound.
